\scanpage{3}
% archive sheet number: 2

\[
H^{\unsure{i|j}}(X^T,F^T)=0,\quad H^0(X^T,F^T)\simeq H^0(X,F)
\]

\deleted{Pour} La \unsure{deuxième} relation triviale par définition de $F^T$.

Pour la première, \deleted{\unsure{on est}} (à prouver pour \uline{tt} \unsure{ouvert} affine) on est ramené par suite spectrale de Cartan : prouver que pour tt \deleted{\unsure{recouvrement} de $X$} famille couvrante $\unsure{\mathfrak X|\mathcal X}=(X_i\longrightarrow X)$, finie fid plate, $\check H^p(\mathfrak X,F)=0$ pour tt $p\ne0$.

\uline{Or} comme $F$ transforme sommes en produits, on peut supposer $\mathfrak X$ \unsure{réduit|réduite} à \uline{un} morphisme fid plat \deleted{\unsure{qu. cpct|quasi-compact}} de schémas affines. On est alors sous les conditions du Th. 1.

Utilisant la suite spectrale de Leray, on trouve

\uline{Corollaire 4} Soit $S'\longrightarrow S$ couvrant (pour la top. fpqc), $F$ qu. coh sur $S$. Alors on a une suite spectrale

\[
H^*(S,F)\Longleftarrow \check H^p(S'/S,{\unsure{\mathscr H|\mathcal H}}^q(F))
\]

[où $\mathscr H^q(F)$ désigne le foncteur $T\longmapsto H^q(T,F\otimes_S T)$ sur $(Sch)_{/S}^{\circ}$]
